% Bagian: turing-machines
% Bab: machines-computations
% Subbagian: disciplined-machines

\documentclass[../../../include/open-logic-section]{subfiles}

\begin{document}

\olfileid[id]{tur}{mac}{dis}
\olsection{Mesin Disiplin}

\begin{explain}
Dalam bagian \olref[hal]{sec}, kita membahas mesin Turing yang memiliki satu
keadaan penghentian khusus~$h$---jika mesin seperti itu berhenti, mesin
dipastikan berhenti dalam keadaan~$h$. Dengan demikian, mesin yang hanya
memiliki satu keadaan penghentian lebih ``disiplin'' daripada mesin Turing
pada umumnya yang kita izinkan. Ada pembatasan lain yang dapat kita kenakan
pada perilaku mesin Turing. Misalnya, kita juga belum melarang mesin Turing
menghapus penanda ujung pita pada petak~$0$, ataupun mencoba bergerak ke kiri
dari petak~$0$. (Definisi kita menyatakan bahwa dalam kasus ini kepala hanya
diam di petak~$0$; definisi lain membuat mesin berhenti.) Kadang-kadang kita
juga perlu dapat mengandaikan bahwa jika suatu mesin Turing berhenti, mesin
itu berhenti di petak~$1$.
\end{explain}

\begin{defn}\ollabel{defn:disciplined}
Suatu mesin Turing~$M$ disebut \emph{disiplin} jika dan hanya jika
\begin{enumerate}
    \item mesin memiliki tepat satu keadaan penghentian khusus~$h$,
    \item jika mesin berhenti, mesin berhenti ketika memindai petak~$1$,
    \item mesin tidak pernah menghapus simbol~$\TMendtape$ pada petak~$0$, dan
    \item mesin tidak pernah mencoba bergerak ke kiri dari petak~$0$.
\end{enumerate}
\end{defn}

\begin{explain}
Kita telah membahas bahwa setiap mesin Turing dapat diubah menjadi mesin
dengan perilaku yang sama tetapi memiliki keadaan penghentian khusus. Caranya
hanyalah dengan menambahkan keadaan baru~$h$, lalu menambahkan instruksi
$\delta(q, \sigma) = \tuple{h, \sigma, N}$ untuk setiap pasangan
$\tuple{q,\sigma}$ tempat $\delta$ semula tidak terdefinisi. Meskipun
pembuktiannya menjemukan, benar bahwa setiap mesin Turing~$M$ dapat diubah
menjadi mesin Turing disiplin~$M'$ yang berhenti pada masukan yang sama dan
menghasilkan keluaran yang sama. Misalnya, jika mesin Turing berhenti tetapi
tidak berada di petak~$1$, kita dapat menambahkan beberapa instruksi agar
kepalanya bergerak ke kiri sampai menemukan penanda ujung pita, lalu bergerak
satu petak ke kanan, kemudian berhenti. Kita serahkan kepada Anda untuk
memikirkan cara menangani syarat-syarat yang lain.
\end{explain}

\begin{ex}
    Dalam \olref{fig:adder-disc}, mesin penjumlahan dari
    \olref[una]{ex:adder} kita ubah menjadi mesin disiplin.
    \begin{figure}\[
\begin{tikzpicture}[->,>=stealth',shorten >=1pt,auto,node distance=2.8cm,
                    semithick]
  \tikzstyle{every state}=[fill=none,draw=black,text=black]

  \node[initial,state]         (A)              {$q_0$};
  \node[state]         (B) [right of=A] {$q_1$};
  \node[state]         (C) [below right of=B] {$q_2$};
  \node[state]         (D) [below left of=C] {$q_3$};
  \node[state]         (H) [left of=D] {$h$};

  \path (A) edge node {\TMtrans{\TMblank}{\TMstroke}{\TMstay}} (B)
           edge [loop below] node {\TMtrans{\TMstroke}{\TMstroke}{\TMright}} (A)
        (B) edge [loop below] node {\TMtrans{\TMstroke}{\TMstroke}{\TMright}} (B)
            edge node {\TMtrans{\TMblank}{\TMblank}{\TMleft}} (C)
        (C) edge node {\TMtrans{\TMstroke}{\TMblank}{\TMleft}} (D)
        (D) edge [loop above] node {\TMtrans{\TMstroke}{\TMstroke}{\TMleft}} (D)
        edge node {\TMtrans{\TMendtape}{\TMendtape}{\TMright}} (H);
\end{tikzpicture}
\]\caption{Mesin penjumlahan disiplin}
\ollabel{fig:adder-disc}
\end{figure}
\end{ex}

\begin{prop}\ollabel{prop:disciplined} Untuk setiap mesin Turing~$M$,
terdapat mesin Turing disiplin~$M'$ yang berhenti dengan keluaran~$O$ jika
$M$ berhenti dengan keluaran~$O$, dan tidak berhenti jika $M$ tidak berhenti.
Secara khusus, setiap fungsi $f\colon\Nat^n \to \Nat$ yang dapat dihitung oleh
mesin Turing juga dapat dihitung oleh mesin Turing disiplin.
\end{prop}

\begin{prob}\label{tur:mac:dis:prob:disc-succ}
Berikan suatu mesin disiplin yang menghitung $f(x) = x+1$.
\end{prob}


\begin{prob}\label{tur:mac:dis:prob:copier}
Carilah mesin disiplin yang, ketika dijalankan pada masukan $\TMstroke^n$,
menghasilkan keluaran $\TMstroke^n \concat \TMblank \concat \TMstroke^n$.
\end{prob}

\end{document}
